% Standalone section fragment for inclusion in an acmart paper.
% Required packages/environments: amsmath, amssymb, amsthm, and algorithm2e.
% The paper preamble must define theorem, lemma, and definition environments.

\section{Formal Foundations of Universal Counting and Replay}
\label{sec:formal-foundations}

Let $G=(V,E)$ be a finite simple undirected graph.
%
For $t\geq 1$, let $\mathcal K_t(G)$ denote the set of $t$-vertex cliques of $G$.
%
We fix integers $1\leq S_{\min}\leq S_{\max}$ and prove that one class-based succinct clique tree for this interval supports exact counting and peeling for every admissible pair $(r,s)$.

\subsection{Universal class-CPI counting}
\label{subsec:universal-class-cpi}

Let $\mathcal M$ be the set of maximal cliques of $G$ having size at least $S_{\min}$.
%
For every vertex $v$, define its relevant-region profile by
\[
\rho(v)=\{M\in\mathcal M:v\in M\}.
\]
Vertices with the same nonempty profile form an \emph{active class}, while every vertex with empty profile forms its own singleton inactive class.
%
Let $\mathcal C$ be the resulting partition of $V$, let $c\subseteq V$ also denote the vertex set of class $c$, and let $w_c=|c|$.
%
The relevant quotient $Q$ has node set $\mathcal C$ and an edge $cd$ exactly when $\rho(c)\cap\rho(d)\neq\emptyset$.

\begin{lemma}[Relevant-class symmetry]
\label{lem:relevant-class-symmetry}
The following properties hold.
\begin{enumerate}
\item Every class $c$ induces a clique, and every quotient edge $cd$ is a complete join between $c$ and $d$.
\item Every clique $X\in\mathcal K_t(G)$ with $t\geq S_{\min}$ uses only active classes, and the classes met by $X$ induce a clique of $Q$.
\item Conversely, choosing an arbitrary number of vertices from each node of a quotient clique always gives a clique of $G$.
\item Every permutation of vertices inside each class maps $\mathcal K_t(G)$ to itself for every $t\in[S_{\min},S_{\max}]$.
\end{enumerate}
The last property concerns the relevant clique families; it does not assert that such a permutation preserves edges that occur only in maximal cliques smaller than $S_{\min}$.
\end{lemma}
\begin{proof}
An inactive class is a singleton and hence induces a clique.
%
For an active class $c$, choose $M\in\rho(c)$.
%
Every vertex of $c$ belongs to $M$, so this class also induces a clique.
%
If $cd\in E(Q)$, choose $M\in\rho(c)\cap\rho(d)$.
%
Then $c\cup d\subseteq M$, so every vertex of $c$ is adjacent to every vertex of $d$.

Now let $X\in\mathcal K_t(G)$ with $t\geq S_{\min}$.
%
Extend $X$ to a maximal clique $M$.
%
Since $|M|\geq t\geq S_{\min}$, we have $M\in\mathcal M$.
%
Thus every vertex of $X$ has a nonempty profile, and any two classes met by $X$ both contain $M$ in their profiles and hence are adjacent in $Q$.
%
Conversely, if $D\subseteq\mathcal C$ is a quotient clique, then every class in $D$ is internally complete and every two distinct classes in $D$ are completely joined by the first paragraph.
%
Therefore every vertex subset of $\bigcup_{c\in D}c$ is a clique.

Finally, let $\gamma$ permute vertices independently inside the classes and let $X\in\mathcal K_t(G)$ for $t\geq S_{\min}$.
%
Choose a relevant maximal clique $M\supseteq X$.
%
If $v\in X\cap c$, then $M\in\rho(c)$, so every vertex of $c$, including $\gamma(v)$, belongs to $M$.
%
Consequently $\gamma(X)\subseteq M$ and is a $t$-clique.
%
Applying the same argument to $\gamma^{-1}$ proves that $\gamma$ maps $\mathcal K_t(G)$ bijectively to itself.
\end{proof}

We use the convention $\binom ab=0$ unless $0\leq b\leq a$ and $a,b$ are integers.

\begin{definition}[Class-SCT leaf and exact interval cover]
\label{def:class-sct-leaf}
A \emph{class-SCT leaf} $L$ specifies, for every $c\in\mathcal C$, an available set $V_{L,c}\subseteq c$, its capacity $n_{L,c}=|V_{L,c}|$, and integer bounds
\[
0\leq \ell_{L,c}\leq u_{L,c}\leq n_{L,c}.
\]
Only finitely many capacities are nonzero, and $\bigcup_cV_{L,c}$ is required to be a clique.
%
For an integer $t$, the $t$-slice encoded by $L$ is
\[
\mathcal K_t(L)=
\left\{
X\subseteq\bigcup_cV_{L,c}:
\substack{
|X|=t,\\
\ell_{L,c}\leq |X\cap c|\leq u_{L,c}\ \text{for every }c}
\right\}.
\]
A finite leaf family $\mathcal L$ is an \emph{exact class-SCT for $[S_{\min},S_{\max}]$} when, for every $t$ in that interval,
\[
\mathcal K_t(G)=\biguplus_{L\in\mathcal L}\mathcal K_t(L),
\]
where $\biguplus$ denotes disjoint union.
%
A leaf is \emph{class-closed} when each $V_{L,c}$ is either $c$ or $\emptyset$.
\end{definition}

Fix an $r$-clique $R$ and write $b_c=|R\cap c|$, so $\sum_cb_c=r$.
%
Its class-composition pattern is $P=(b_c)_{c\in\mathcal C}$.
%
We say that $L$ \emph{hosts $R$} when $R\cap c\subseteq V_{L,c}$ for every $c$.
%
For a class-closed leaf, hosting depends only on $P$, and we then say that $L$ hosts $P$.

\begin{lemma}[General class-leaf count]
\label{lem:general-class-leaf-count}
If $L$ hosts $R$, then the number of $s$-cliques in the leaf slice that contain $R$ is
\begin{equation}
\label{eq:general-class-leaf-count}
\operatorname{count}_{L,s}(P)
=
\sum_{\substack{
(y_c)_c\in\mathbb Z_{\geq0}^{\mathcal C}:\ \sum_cy_c=s\\
\max\{\ell_{L,c},b_c\}\leq y_c\leq u_{L,c}\ \forall c}}
\prod_{c\in\mathcal C}
\binom{n_{L,c}-b_c}{y_c-b_c}.
\end{equation}
If $L$ does not host $R$, define $\operatorname{count}_{L,s}(P)=0$.
\end{lemma}
\begin{proof}
Fix a feasible composition $\mathbf y=(y_c)_c$ in Equation~\eqref{eq:general-class-leaf-count}.
%
The desired clique already contains the $b_c$ vertices of $R\cap c$ and must choose exactly $y_c-b_c$ additional vertices from the $n_{L,c}-b_c$ vertices of $V_{L,c}\setminus R$.
%
The classwise choices are independent, so the number of choices is the displayed product.
%
Every chosen set lies in $\mathcal K_s(L)$ because its class counts obey the leaf bounds, and it is a clique because the union of the available sets is a clique.
%
Conversely, every member of $\mathcal K_s(L)$ containing $R$ has a unique class-count vector $\mathbf y$ and a unique additional subset in every class.
%
Distinct vectors or distinct classwise choices give distinct cliques, so the sum is bijective and has no overcount.
%
If $L$ does not host $R$, no subset of $\bigcup_cV_{L,c}$ can contain $R$, which proves the zero case.
\end{proof}

The common one-pool formula is an exact specialization of this count.
%
After conditioning on $R$, suppose a leaf offers one unrestricted pool of $w$ additional vertices and no other choice.
%
The only free coordinate chooses $s-r$ vertices, and Equation~\eqref{eq:general-class-leaf-count} becomes
\begin{equation}
\label{eq:single-pool-special-case}
\operatorname{count}_{L,s}(P)=\binom{w}{s-r}.
\end{equation}

\begin{theorem}[Universal CPI, independent of $r$]
\label{thm:universal-cpi}
Let $\mathcal L$ be one exact class-closed class-SCT built for $[S_{\min},S_{\max}]$, without taking $r$ as a construction parameter.
%
For every $r\geq1$ with $r+1\leq S_{\max}$, every $r$-clique $R$ with pattern $P$, and every
\[
s\in[\max\{r+1,S_{\min}\},S_{\max}],
\]
the exact initial $s$-support of $R$ is
\begin{equation}
\label{eq:universal-support}
\operatorname{sup}_s(R)
:=
|\{S\in\mathcal K_s(G):R\subseteq S\}|
=
\sum_{L\in\mathcal L}\operatorname{count}_{L,s}(P).
\end{equation}
Thus one immutable class-SCT is an exact counting representation for every cell in the bounded $(r,s)$ plane; $r$ enters only through $\mathbf b$ and $s$ only through the constraint $\sum_cy_c=s$.
\end{theorem}
\begin{proof}
Fix $r$, $s$, and $R$.
%
The exact-cover property in Definition~\ref{def:class-sct-leaf} partitions $\mathcal K_s(G)$ into the disjoint leaf slices $\mathcal K_s(L)$.
%
Intersecting each slice with $\{S:R\subseteq S\}$ preserves both coverage and disjointness.
%
Lemma~\ref{lem:general-class-leaf-count} gives the cardinality of each resulting part, so additivity of cardinality over a disjoint union proves Equation~\eqref{eq:universal-support}.
%
If $R$ belongs to no $s$-clique, every intersection is empty and both sides are zero.
%
Neither the leaf definition nor the exact-cover property contains $r$, while the same leaf bounds define every $s$-slice in the construction interval.
%
Therefore changing the query order or target size does not rebuild or mutate $\mathcal L$.
\end{proof}

\subsection{Construction by a disjoint pivot recurrence}
\label{subsec:gen-construction}

We now give a self-contained weighted-class version of the canonical degeneracy and pivot partition used by Pivoter~\cite{Pivoter}.
%
Let $\pi$ be a fixed degeneracy order of the quotient nodes, and write $N_Q^+(c)=\{d\in N_Q(c):\pi(d)>\pi(c)\}$.
%
A weighted quotient clique is a vector $\mathbf x$ with $0\leq x_c\leq w_c$ whose positive coordinates induce a clique of $Q$; it represents every concrete choice of $x_c$ vertices from each class.

\begin{definition}[State and completion of \textsc{Gen}]
\label{def:gen-state}
A \textsc{Gen} state is $(F,U,A,[S_{\min},S_{\max}])$.
%
The \emph{spine} $F$ is a map with a fixed positive multiplicity $f_c\in[1,w_c]$ on every $c\in\operatorname{dom}(F)$, the \emph{free pool} $U$ is a set of classes with multiplicity range $[0,w_c]$, and the \emph{open pool} $A$ contains undecided classes.
%
The three class sets are disjoint, $\operatorname{dom}(F)\cup U$ induces a quotient clique, and every class in $A$ is adjacent to every class in $\operatorname{dom}(F)\cup U$.
%
A concrete completion is a clique $X$ whose class counts $x_c=|X\cap c|$ satisfy
\[
\begin{aligned}
x_c&=f_c &&(c\in\operatorname{dom}(F)),\\
0\leq x_c&\leq w_c &&(c\in U\cup A),\\
x_c&=0 &&(c\notin\operatorname{dom}(F)\cup U\cup A),
\end{aligned}
\]
whose positive coordinates induce a quotient clique, and whose size lies in $[S_{\min},S_{\max}]$.
\end{definition}

For a state, let $\mathfrak C(F,U,A)$ denote its concrete completion family.
%
Choose a deterministic pivot $p\in A$, and let $B=A\setminus(\{p\}\cup N_Q(p))=\{v_1<\cdots<v_q\}$ under a fixed class order.
%
If $B=\emptyset$, the recurrence is
\begin{equation}
\label{eq:gen-universal-recurrence}
\mathfrak C(F,U,A)
=
\mathfrak C(F,U\cup\{p\},A\setminus\{p\}).
\end{equation}
If $B\neq\emptyset$, define
\[
A_0=A\setminus B,
\qquad
A_t=(A\cap N_Q(v_t))\setminus\{v_1,\ldots,v_t\},
\]
and obtain the disjoint recurrence
\begin{equation}
\label{eq:gen-nonuniversal-recurrence}
\mathfrak C(F,U,A)
=
\mathfrak C(F,U,A_0)
\ \dot\cup\!
\mathop{\dot\bigcup}_{t=1}^{q}
\mathop{\dot\bigcup}_{j=1}^{w_{v_t}}
\mathfrak C(F\cup\{v_t\mapsto j\},U,A_t),
\end{equation}
where a child with forced size above $S_{\max}$ is empty.
%
The first family uses no class in $B$, while a $(t,j)$ family uses $v_t$ with multiplicity $j$ and uses no earlier member of $B$.

\begin{algorithm}[t]
\caption{\textsc{BuildClassSCT} and its recursive procedure \textsc{Gen}.}
\label{alg:gen-class-sct}
\DontPrintSemicolon
\SetKwFunction{Gen}{Gen}
\SetKwProg{Fn}{Function}{:}{end}
\KwIn{quotient $Q$, class weights $(w_c)_c$, degeneracy order $\pi$, interval $[S_{\min},S_{\max}]$}
\KwOut{leaf family $\mathcal L$}
$\mathcal L\gets\emptyset$
\;
\Fn{\Gen{$F,U,A$}}{
  $a\gets\sum_{c\in\operatorname{dom}(F)}f_c$
  \;
  \If{$a>S_{\max}$ \textbf{or} $a+\sum_{c\in U\cup A}w_c<S_{\min}$}{\Return}
  \If{$A=\emptyset$}{
    create class-closed $L$ with $V_{L,c}=c$ on $\operatorname{dom}(F)\cup U$ and $V_{L,c}=\emptyset$ otherwise
    \;
    set $(n_{L,c},\ell_{L,c},u_{L,c})\gets(w_c,f_c,f_c)$ for $c\in\operatorname{dom}(F)$
    \;
    set $(n_{L,c},\ell_{L,c},u_{L,c})\gets(w_c,0,w_c)$ for $c\in U$ and $(0,0,0)$ otherwise
    \;
    $\mathcal L\gets\mathcal L\cup\{L\}$; \Return
    \;
  }
  choose the deterministic pivot $p\in A$; $B\gets A\setminus(\{p\}\cup N_Q(p))$
  \;
  \eIf{$B=\emptyset$}{
    \Gen{$F,U\cup\{p\},A\setminus\{p\}$}
    \;
  }{
    order $B$ as $v_1<\cdots<v_q$; \Gen{$F,U,A\setminus B$}
    \;
    \For{$t\gets1$ \KwTo $q$}{
      $A_t\gets(A\cap N_Q(v_t))\setminus\{v_1,\ldots,v_t\}$
      \;
      \For{$j\gets1$ \KwTo $w_{v_t}$}{
        \Gen{$F\cup\{v_t\mapsto j\},U,A_t$}
        \;
      }
    }
  }
}
\ForEach{$c\in\mathcal C$ in increasing $\pi$ order}{
  \For{$j\gets1$ \KwTo $\min\{w_c,S_{\max}\}$}{
    \Gen{$\{c\mapsto j\},\emptyset,N_Q^+(c)$}
    \;
  }
}
\Return{$\mathcal L$}
\;
\end{algorithm}

The pivot may, for example, maximize $\sum_{z\in A\cap N_Q(p)}w_z$; the choice affects the tree size but not the recurrence's correctness.

\begin{theorem}[Disjoint leaves and exact cover]
\label{thm:gen-exact-cover}
For every valid call $\textsc{Gen}(F,U,A)$, the leaves emitted below that call form a pairwise-disjoint partition of $\mathfrak C(F,U,A)$.
%
Across the outer degeneracy seeds in Algorithm~\ref{alg:gen-class-sct}, every concrete clique of every size in $[S_{\min},S_{\max}]$ belongs to exactly one emitted leaf slice.
%
Consequently Algorithm~\ref{alg:gen-class-sct} constructs one exact class-SCT for the entire interval, independently of $r$.
\end{theorem}
\begin{proof}
We first prove the per-call assertion by induction on $|A|$.
%
If either pruning test succeeds, the forced multiplicity already exceeds $S_{\max}$ or even using every remaining vertex cannot reach $S_{\min}$, so $\mathfrak C(F,U,A)=\emptyset$ and emitting no leaf is exact.
%
If $A=\emptyset$ and neither pruning test succeeds, the emitted leaf chooses exactly $f_c$ arbitrary vertices in every spine class and an arbitrary number from $0$ through $w_c$ in every free class.
%
These choices are precisely the completions in the size window, so the single leaf is an exact one-part cover.

Assume $A\neq\emptyset$ and that the assertion holds for smaller open pools.
%
If $B=\emptyset$, the pivot is adjacent to every other open class and, by the state invariant, to every spine and free class.
%
Its multiplicity can therefore vary independently from $0$ through $w_p$.
%
Moving $p$ from $A$ to $U$ records exactly that choice, so Equation~\eqref{eq:gen-universal-recurrence} is an equality of concrete completion families.
%
The child has a smaller open pool, and the induction hypothesis applies.

Now suppose $B=\{v_1,\ldots,v_q\}\neq\emptyset$ and take any $X\in\mathfrak C(F,U,A)$.
%
If $X$ uses no class in $B$, then every open class used by $X$ lies in $A_0=A\setminus B$, so $X$ belongs to the first child of Equation~\eqref{eq:gen-nonuniversal-recurrence}.
%
Conversely, no completion of that child uses a class in $B$.
%
Otherwise let $v_t$ be the first class of $B$ used by $X$ and let $j=|X\cap v_t|$.
%
Because $X$ is a clique, every other open class used by $X$ is adjacent to $v_t$.
%
The pivot $p$ is not used because $pv_t\notin E(Q)$, and no $v_i$ with $i<t$ is used by the choice of $t$.
%
Thus every remaining open class used by $X$ lies in $A_t$, and $X$ belongs to the unique $(t,j)$ child.
%
Conversely, the state invariant and the definition of $A_t$ ensure that every completion of the $(t,j)$ child is a clique, uses $v_t$ exactly $j$ times, and uses no earlier member of $B$.
%
Hence the branch labels ``no member of $B$'' and ``first used member $v_t$ with multiplicity $j$'' are mutually exclusive and exhaustive.
%
This proves that Equation~\eqref{eq:gen-nonuniversal-recurrence} is a disjoint union.
%
Every child has a strictly smaller open pool, so the induction hypothesis partitions each child into leaves, and their union proves the per-call claim.

It remains to prove disjointness across seeds.
%
Let $X$ be a nonempty concrete clique in the construction window, let $c$ be the unique $\pi$-minimum class met by $X$, and let $j=|X\cap c|$.
%
Every other class met by $X$ is a later quotient neighbour of $c$, so $X$ is a completion of the seed $(\{c\mapsto j\},\emptyset,N_Q^+(c))$.
%
No other seed can contain $X$, because its minimum used class or its multiplicity at that class would differ.
%
The per-call partition therefore places $X$ in exactly one emitted leaf.
%
The argument holds simultaneously for all sizes in the interval, proving the exact-cover property in Definition~\ref{def:class-sct-leaf}.
%
Algorithm~\ref{alg:gen-class-sct} has no order-$r$ input, which proves the final assertion.
\end{proof}

\begin{theorem}[Additive support]
\label{thm:additive-support}
Let $\mathcal L$ be the leaves emitted by Algorithm~\ref{alg:gen-class-sct}, fix $s\in[S_{\min},S_{\max}]$, and let $R$ have pattern $P$.
%
Then
\begin{equation}
\label{eq:additive-support}
\operatorname{sup}_s(R)
=
\sum_{L:\,L\text{ hosts }P}\operatorname{count}_{L,s}(P),
\end{equation}
with no inclusion--exclusion between leaves.
%
More generally, for any set $Z\subseteq\mathcal K_s(G)$ of dead witnesses, if
\[
\operatorname{count}^{Z}_{L,s}(P)
=
|\{S\in\mathcal K_s(L)\setminus Z:R\subseteq S\}|,
\]
then the surviving support is the plain sum $\sum_L\operatorname{count}^{Z}_{L,s}(P)$.
\end{theorem}
\begin{proof}
Theorem~\ref{thm:gen-exact-cover} gives the disjoint partition
\[
\mathcal K_s(G)=\biguplus_L\mathcal K_s(L).
\]
%
After intersecting every part with the witnesses containing $R$, nonhosting leaves contribute the empty set and hosting leaves contribute the sets counted by Lemma~\ref{lem:general-class-leaf-count}.
%
Cardinality is additive over this disjoint union, which proves Equation~\eqref{eq:additive-support}.
%
For the second statement, the families $\mathcal K_s(L)\setminus Z$ remain pairwise disjoint and cover $\mathcal K_s(G)\setminus Z$.
%
Intersecting them with the witnesses containing $R$ and adding their cardinalities proves the surviving formula.
%
Inclusion--exclusion may be used internally to evaluate the union of several deletion conditions inside one leaf, but no cross-leaf correction exists or is needed.
\end{proof}

\subsection{Weighted pattern orbits and exact peeling}
\label{subsec:weighted-pattern-peeling}

Fix a cell $(r,s)$ with $1\leq r<s$ and $s\in[S_{\min},S_{\max}]$.
%
Let $\mathcal W=\mathcal K_s(G)$ be the witness set and let
\[
\mathcal U=\bigcup_{S\in\mathcal W}\mathcal K_r(S)
\]
be the individual $r$-clique peeling units that occur in at least one witness.
%
Every $r$-clique outside $\mathcal U$ has support and core value zero and may be returned separately.
%
For $A\subseteq\mathcal U$, define
\begin{equation}
\label{eq:induced-witness-support}
\begin{aligned}
\mathcal W(A)&=\{S\in\mathcal W:\mathcal K_r(S)\subseteq A\},\\
d_A(R)&=|\{S\in\mathcal W(A):R\subseteq S\}|.
\end{aligned}
\end{equation}
Thus a witness survives exactly while all of its $r$-subcliques survive.

For an integer $k\geq0$, call $A$ \emph{$k$-feasible} when $d_A(R)\geq k$ for every $R\in A$.
%
The union of two $k$-feasible sets is $k$-feasible because $d_A(R)$ is nondecreasing in $A$.
%
Hence there is a unique largest $k$-feasible set, denoted $K_k$, and we define
\begin{equation}
\label{eq:core-threshold-definition}
\kappa_{r,s}(R)=\max\{k:R\in K_k\}.
\end{equation}

\begin{lemma}[Threshold characterization of peeling]
\label{lem:threshold-peeling}
In any finite unit--witness system with supports monotone in the survivor set, repeatedly deleting any unit whose current support is below $k$ terminates at the unique largest $k$-feasible set $K_k$, independently of deletion order.
%
Consequently, ordinary minimum-support peeling with the clamp $\lambda\gets\max\{\lambda,d_A(x)\}$ assigns every unit $x$ the value $\max\{k:x\in K_k\}$.
\end{lemma}
\begin{proof}
During threshold-$k$ deletion, no member of $K_k$ can be removed.
%
Indeed, if the current set $A$ contains $K_k$ and $x\in K_k$, monotonicity gives $d_A(x)\geq d_{K_k}(x)\geq k$.
%
The terminal set $A^*$ is itself $k$-feasible because no remaining unit has support below $k$.
%
Since $K_k$ survives throughout and is the largest $k$-feasible set, $K_k\subseteq A^*\subseteq K_k$, so $A^*=K_k$.

Run minimum-support peeling and fix $k$.
%
Whenever the current survivor set strictly contains $K_k$, it is not $k$-feasible by maximality, so its minimum support is below $k$ and the next minimum-support deletion is also a valid threshold-$k$ deletion.
%
Therefore all units outside $K_k$ disappear before the clamped level reaches $k$.
%
When the survivor set first becomes $K_k$, every current support is at least $k$, so the next deletion raises the clamped level to at least $k$ and the clamp keeps all later assigned values at least $k$.
%
Thus a unit receives value at least $k$ exactly when it belongs to $K_k$.
%
This equivalence for every integer $k$ proves the claimed value.
\end{proof}

A valid pattern $P$ is a composition $\mathbf b^P=(b_c^P)_c$ with $\sum_cb_c^P=r$ for which
\[
\mathcal U_P=\{R\in\mathcal U:|R\cap c|=b_c^P\ \text{for every }c\}
\]
is nonempty.
%
Lemma~\ref{lem:relevant-class-symmetry} implies that arbitrary within-class permutations act transitively on $\mathcal U_P$ and preserve the unit--witness incidence relation.
%
Therefore $\mathcal U_P$ contains every independent choice of $b_c^P$ vertices from each class: one realization lies in an $s$-witness, and applying all within-class permutations to that witness produces every such choice in an $s$-witness.
%
The pattern super-node consequently has multiplicity
\begin{equation}
\label{eq:pattern-multiplicity}
m(P)=|\mathcal U_P|=\prod_{c\in\mathcal C}\binom{w_c}{b_c^P}.
\end{equation}
Its current support is the support of one representative, not the sum over its $m(P)$ representatives.

We next define the exact multiplicity-aware cascade.
%
The class-SCT emitted by Algorithm~\ref{alg:gen-class-sct} is class-closed, so leaf hosting and all following tests depend only on compositions.
%
Let $\mathcal D$ be the set of pattern orbits already deleted.
%
For an alive pattern $Q$ and a leaf $L$, define
\begin{equation}
\label{eq:live-leaf-count}
\operatorname{LiveCount}_{L,s}^{\mathcal D}(Q)
=
\sum_{\substack{
\sum_cy_c=s\\
\max\{\ell_{L,c},b_c^Q\}\leq y_c\leq u_{L,c}\ \forall c\\
\mathbf b^D\nleq\mathbf y\ \forall D\in\mathcal D}}
\prod_c\binom{n_{L,c}-b_c^Q}{y_c-b_c^Q},
\end{equation}
with value zero when $L$ does not host $Q$.
%
Here $\mathbf a\leq\mathbf y$ means componentwise inequality.
%
An $s$-clique of composition $\mathbf y$ contains some representative of pattern $D$ exactly when $\mathbf b^D\leq\mathbf y$.
%
Thus
\begin{equation}
\label{eq:pattern-live-support}
\sigma_{\mathcal D}(Q)=\sum_{L\in\mathcal L}\operatorname{LiveCount}_{L,s}^{\mathcal D}(Q)
\end{equation}
is the exact current per-representative support.

When an alive pattern $P$ is deleted, the exact decrement to a distinct alive pattern $Q$ is
\begin{equation}
\label{eq:aggregate-pattern-decrement}
\Delta_{\mathcal D}(P\to Q)
=
\sum_{L\in\mathcal L}
\sum_{\substack{
\sum_cy_c=s\\
\max\{\ell_{L,c},b_c^P,b_c^Q\}\leq y_c\leq u_{L,c}\ \forall c\\
\mathbf b^D\nleq\mathbf y\ \forall D\in\mathcal D}}
\prod_c\binom{n_{L,c}-b_c^Q}{y_c-b_c^Q}.
\end{equation}
The cascade rule is
\begin{equation}
\label{eq:aggregate-pattern-cascade}
\sigma_{\mathcal D\cup\{P\}}(Q)
=
\sigma_{\mathcal D}(Q)-\Delta_{\mathcal D}(P\to Q).
\end{equation}

\begin{theorem}[Weighted-pattern peeling equivalence]
\label{thm:weighted-pattern-equivalence}
Run minimum-support peeling on pattern super-nodes.
%
When pattern $P$ is removed, assign the clamped peel value to all $m(P)$ representatives, delete the entire orbit $\mathcal U_P$, and update every other alive pattern by Equations~\eqref{eq:aggregate-pattern-decrement}--\eqref{eq:aggregate-pattern-cascade}.
%
Then every representative $R\in\mathcal U_P$ receives exactly its individual from-scratch value $\kappa_{r,s}(R)$.
%
The multiplicity $m(P)$ is used for reporting and for the number of units removed, but it is not a multiplier in either support or decrement.
\end{theorem}
\begin{proof}
Let $\Gamma=\prod_{c\in\mathcal C}\operatorname{Sym}(c)$ act by independent within-class permutations.
%
By Lemma~\ref{lem:relevant-class-symmetry}, $\Gamma$ maps $s$-cliques to $s$-cliques and preserves containment, so it is an automorphism group of the unit--witness system in Equation~\eqref{eq:induced-witness-support}.
%
Its orbits on $\mathcal U$ are exactly the pattern realization sets: the action preserves every class count, and symmetric groups are transitive on equal-size subsets of each class.

The largest individual $k$-feasible set $K_k$ is unique.
%
For any $\gamma\in\Gamma$, the set $\gamma(K_k)$ is also $k$-feasible, so uniqueness gives $\gamma(K_k)=K_k$.
%
Hence $K_k$ is a union of complete pattern orbits.
%
For a pattern set $\mathcal B$, let $A(\mathcal B)=\bigcup_{P\in\mathcal B}\mathcal U_P$.
%
All representatives of a pattern have the same support in this invariant survivor set.
%
Therefore $A(\mathcal B)$ is individually $k$-feasible if and only if every pattern in $\mathcal B$ has per-representative support at least $k$.
%
The unique largest feasible pattern set is consequently
\[
\overline K_k=\{P:\mathcal U_P\subseteq K_k\}.
\]
Lemma~\ref{lem:threshold-peeling} applied to the pattern system shows that minimum-support pattern peeling assigns $P$ the value $\max\{k:P\in\overline K_k\}$.
%
For every $R\in\mathcal U_P$, this equals $\max\{k:R\in K_k\}=\kappa_{r,s}(R)$.

It remains to prove that the aggregate cascade maintains the support used in this quotient argument.
%
Fix a representative $R_Q\in\mathcal U_Q$, a leaf $L$, and a witness $S\in\mathcal K_s(L)$ containing $R_Q$ with class composition $\mathbf y$.
%
The witness contains a representative of pattern $P$ if and only if $b_c^P\leq y_c$ for every class: necessity follows from class counts, and sufficiency follows by choosing any $b_c^P$ of the $y_c$ witness vertices in each class.
%
The same criterion for every $D\in\mathcal D$ says that $S$ is live before deleting $P$ exactly when $\mathbf b^D\nleq\mathbf y$ for all $D$.
%
For fixed $L$ and $\mathbf y$, Lemma~\ref{lem:general-class-leaf-count} counts the witnesses containing $R_Q$ by the product in Equation~\eqref{eq:aggregate-pattern-decrement}.
%
Thus that equation counts exactly the witnesses containing $R_Q$ that are live before the deletion and die because they contain at least one representative of $P$.
%
Theorem~\ref{thm:additive-support} makes the sum over leaves disjoint, so Equation~\eqref{eq:aggregate-pattern-cascade} is exact.
%
A witness that contains several representatives of $P$ still dies once, and the componentwise predicate counts it once; multiplying by $m(P)$ would therefore overcount.
\end{proof}

\subsection{Shell-order replay and the complete replay algorithm}
\label{subsec:shell-order-replay}

We state replay first for an arbitrary finite unit--witness system and then instantiate it with weighted pattern orbits.
%
Write $\kappa(x)$ for the from-scratch core value of a unit and $K_k=\{x:\kappa(x)\geq k\}$ for its unique largest $k$-feasible set.

\begin{theorem}[Shell-order replay]
\label{thm:shell-order-replay}
Partition the units into a certified set $\mathit{Cert}$ and a residue set $\mathit{Res}$, and suppose the exact value $\kappa(x)$ is known for every $x\in\mathit{Cert}$.
%
Delete certified units in nondecreasing order of their known values, deleting each certified $x$ at level $\kappa(x)$ even when its current support exceeds $\kappa(x)$.
%
Interleave these scheduled deaths with ordinary minimum-current-support peeling of residue units, using the usual clamp to the current level.
%
Before advancing beyond an integer level $k$, process every certified death scheduled at $k$ and every residue cascade whose clamped key is at most $k$; ties at level $k$ may be processed in any order.
%
Then every residue unit is deleted at, and is assigned, its from-scratch value.
%
The conclusion applies to weighted pattern orbits when current supports are per representative and deletions use Equation~\eqref{eq:aggregate-pattern-cascade}.
\end{theorem}
\begin{proof}
We prove by induction on $k\geq0$ that the alive set at the beginning of level $k$ is exactly $K_k$ and that the alive set after level $k$ is exactly $K_{k+1}$.
%
The base case holds because supports are nonnegative, so every unit lies in $K_0$.

Assume that level $k$ begins with alive set $K_k$.
%
Every scheduled death at this level removes a certified unit of exact core value $k$, hence a member of $K_k\setminus K_{k+1}$.
%
No residue unit in $K_{k+1}$ can be removed ordinarily at level $k$.
%
Indeed, throughout level $k$ the current alive set contains $K_{k+1}$ until this claim is applied, and every $x\in K_{k+1}$ has support at least $k+1$ inside $K_{k+1}$.
%
Support monotonicity gives current support at least $k+1$, so its clamped residue key exceeds $k$.
%
By induction over the deletions made within level $k$, every deleted unit therefore belongs to $K_k\setminus K_{k+1}$, and $K_{k+1}$ remains alive.

We show that level $k$ cannot terminate while a shell unit in $K_k\setminus K_{k+1}$ remains.
%
After all certified deaths scheduled at $k$ have been processed, every remaining certified unit has core at least $k+1$ and hence lies in $K_{k+1}$.
%
Suppose a shell unit remains but no residue unit has current support at most $k$.
%
Every remaining residue unit then has support at least $k+1$.
%
Every remaining certified unit also has support at least $k+1$, because it lies in $K_{k+1}$ and the alive set contains $K_{k+1}$.
%
The entire current alive set would be $(k+1)$-feasible, contradicting the maximality of $K_{k+1}$ because that set contains a shell unit outside $K_{k+1}$.
%
Thus a residue deletion at level $k$ remains available until every shell unit has disappeared.
%
At the end of the level, exactly $K_{k+1}$ remains, completing the induction.

The induction assigns every residue unit its unique shell value $\kappa(x)$.
%
It also explains why a certified core-$k$ unit may be scheduled when its current support is larger than $k$: all support that protects $K_{k+1}$ is internal to $K_{k+1}$ and survives deletion of the core-$k$ shell.
%
For weighted patterns, Theorem~\ref{thm:weighted-pattern-equivalence} identifies pattern superlevel sets with individual superlevel sets, and Equation~\eqref{eq:aggregate-pattern-cascade} supplies the exact monotone current supports.
%
The same induction therefore applies.
\end{proof}

By Theorem~\ref{thm:weighted-pattern-equivalence}, all representatives of a pattern have one common value; write it as $\kappa_{r,s}(P)$.
%
We now define the closed-form certificate used by replay.
%
For an active pattern $P$, let
\[
c(P)=\max\{|M|:M\in\mathcal M,\ R\subseteq M\ \text{for every }R\in\mathcal U_P\}
\]
and define its clique floor at the cell by
\begin{equation}
\label{eq:closed-form-floor}
F_{r,s}(P)=\binom{c(P)-r}{s-r}.
\end{equation}
The maximum exists because an active pattern has a representative in an $s$-witness, every such witness extends to a relevant maximal clique, and that clique contains the entire orbit by class symmetry.
%
The hosting-region set, and hence $c(P)$, is independent of the chosen representative because all members of every used class have the same profile.

\begin{lemma}[Clique floor and valid closed-form certificate]
\label{lem:closed-form-certificate}
Every representative of $P$ has $\kappa_{r,s}$ value at least $F_{r,s}(P)$.
%
If an independently established upper bound $U_{r,s}(P)$ with $\kappa_{r,s}(P)\leq U_{r,s}(P)$ satisfies $U_{r,s}(P)=F_{r,s}(P)$, then
\[
\kappa_{r,s}(P)=F_{r,s}(P),
\]
and $P$ is a valid certified pattern for shell-order replay.
\end{lemma}
\begin{proof}
Choose a hosting maximal clique $M$ of size $c(P)$ and let $A=\{R\in\mathcal K_r(G):R\subseteq M\}$.
%
For every $R\in A$, exactly $\binom{c(P)-r}{s-r}$ witnesses contained in $M$ contain $R$, and all of their $r$-subcliques also belong to $A$.
%
Thus $A$ is $F_{r,s}(P)$-feasible, so every representative of $P$ in $A$ belongs to $K_{F_{r,s}(P)}$ and has core value at least the floor.
%
If the supplied upper bound equals this lower bound, equality follows immediately.
\end{proof}

For a fixed cell, let $\mathcal H_s(P)$ be the leaves for which there exists a feasible vector $\mathbf y$ satisfying the leaf bounds, $\sum_cy_c=s$, and $\mathbf b^P\leq\mathbf y$.
%
These are exactly the leaves that contain an $s$-witness containing a representative of $P$.
%
For every leaf $L$, store a Pareto-minimal antichain $\mathcal A_L$ of deleted pattern vectors.
%
The operation $\operatorname{InsertMin}(\mathcal A_L,\mathbf b)$ discards $\mathbf b$ if some $\mathbf a\in\mathcal A_L$ satisfies $\mathbf a\leq\mathbf b$; otherwise it inserts $\mathbf b$ and removes every $\mathbf a$ with $\mathbf b\leq\mathbf a$.
%
Dominance removal preserves the upward-closed set $\{\mathbf y:\exists\mathbf a\in\mathcal A_L,\mathbf a\leq\mathbf y\}$.
%
Let $\operatorname{LiveCount}(L,s,Q,\mathcal A_L)$ denote Equation~\eqref{eq:live-leaf-count} with the deleted-vector condition represented by this antichain.

An event $(h,\tau,P,v)$ is stale when $P$ is no longer alive, or when $\tau=\mathsf{residue}$ and either $v$ is not the current version of $P$ or $h$ is not its current key.
%
Events with equal key may be popped in any order.

\begin{algorithm*}[t]
\caption{\textsc{ReplayPeel}: exact shell-ordered peeling of certified and residue pattern orbits.}
\label{alg:replay-peel}
\DontPrintSemicolon
\KwIn{cell $(r,s)$; class-SCT $\mathcal L$; active-pattern partition $(\mathit{Cert},\mathit{Res})$; vectors $\mathbf b^P$, multiplicities $m(P)$, and hosting lists $\mathcal H_s(P)$; valid certificates $\widehat\kappa(P)=F_{r,s}(P)=\kappa_{r,s}(P)$ for $P\in\mathit{Cert}$}
\KwOut{$\kappa_{r,s}(Q)$ for every $Q\in\mathit{Res}$, with multiplicity $m(Q)$}
\ForEach{$L\in\mathcal L$}{$\mathcal A_L\gets\emptyset$}
$\mathit{alive}[P]\gets\mathsf{true}$ for every $P\in\mathit{Cert}\cup\mathit{Res}$; initialize an empty min-priority queue $\mathcal Q$
\;
\ForEach{$P\in\mathit{Cert}$}{push $(\widehat\kappa(P),\mathsf{scheduled},P,0)$ into $\mathcal Q$}
\ForEach{$Q\in\mathit{Res}$}{
  $\sigma[Q]\gets\sum_{L\in\mathcal H_s(Q)}\operatorname{LiveCount}(L,s,Q,\mathcal A_L)$
  \;
  $\mathit{key}[Q]\gets\sigma[Q]$; $\mathit{ver}[Q]\gets0$
  \;
  push $(\mathit{key}[Q],\mathsf{residue},Q,\mathit{ver}[Q])$ into $\mathcal Q$
  \;
}
$\lambda\gets0$
\;
\While{$\mathcal Q\neq\emptyset$}{
  $(h,\tau,P,v)\gets\operatorname{popMin}(\mathcal Q)$
  \;
  \If{$\operatorname{Stale}(h,\tau,P,v)$}{\textbf{continue}}
  $\lambda\gets\max\{\lambda,h\}$
  \;
  \eIf{$\tau=\mathsf{scheduled}$}{$\kappa[P]\gets\widehat\kappa(P)$}{$\kappa[P]\gets\lambda$}
  assign $\kappa[P]$ to all $m(P)$ representatives; $\mathit{alive}[P]\gets\mathsf{false}$
  \;
  $\mathit{Affected}\gets\emptyset$
  \;
  \ForEach{$L\in\mathcal H_s(P)$}{
    $\mathcal A_L\gets\operatorname{InsertMin}(\mathcal A_L,\mathbf b^P)$
    \;
    add every alive $Q\in\mathit{Res}$ with $L\in\mathcal H_s(Q)$ to $\mathit{Affected}$
    \;
  }
  \ForEach{$Q\in\mathit{Affected}$}{
    $\sigma[Q]\gets\sum_{L\in\mathcal H_s(Q)}\operatorname{LiveCount}(L,s,Q,\mathcal A_L)$
    \;
    $\mathit{key}[Q]\gets\max\{\lambda,\sigma[Q]\}$; $\mathit{ver}[Q]\gets\mathit{ver}[Q]+1$
    \;
    push $(\mathit{key}[Q],\mathsf{residue},Q,\mathit{ver}[Q])$ into $\mathcal Q$
    \;
  }
}
\Return{$\{(Q,\kappa[Q],m(Q)):Q\in\mathit{Res}\}$}
\;
\end{algorithm*}

\begin{theorem}[Correctness of \textsc{ReplayPeel}]
\label{thm:replay-peel-correctness}
If every supplied certificate is valid, Algorithm~\ref{alg:replay-peel} terminates and returns the from-scratch value $\kappa_{r,s}(Q)$ for every residue pattern $Q$ and hence for all $m(Q)$ individual representatives.
\end{theorem}
\begin{proof}
We first prove the leaf-state invariant: after any deletion prefix, a witness of composition $\mathbf y$ in leaf $L$ is live if and only if no $\mathbf a\in\mathcal A_L$ satisfies $\mathbf a\leq\mathbf y$.
%
Initially every antichain is empty, so the invariant holds.
%
Deleting pattern $P$ kills exactly the previously live witnesses with $\mathbf b^P\leq\mathbf y$, by the containment argument in Theorem~\ref{thm:weighted-pattern-equivalence}.
%
Inserting $\mathbf b^P$ therefore updates the excluded upward-closed set exactly, and deleting dominated antichain vectors does not change that set.
%
The invariant follows by induction on the deletion prefix.
%
Equations~\eqref{eq:live-leaf-count} and~\eqref{eq:pattern-live-support}, together with Theorem~\ref{thm:additive-support}, then show that every recomputed $\sigma[Q]$ is the exact current per-representative support.

If deleting $P$ changes the support of a residue pattern $Q$, some newly dead witness contains representatives of both patterns.
%
That witness lies in one leaf belonging to both $\mathcal H_s(P)$ and $\mathcal H_s(Q)$, so the algorithm includes $Q$ in $\mathit{Affected}$ and recomputes it.
%
Patterns included unnecessarily are merely recomputed to the same exact value.
%
Incrementing versions makes all older residue events stale, so every live residue event accepted by the queue has key $\max\{\lambda,\sigma[Q]\}$ with exact current support.

All certified events are inserted at their exact closed-form core levels, and the min-priority queue processes these levels in nondecreasing order.
%
Residue events implement ordinary minimum-support peeling with the same nondecreasing-level clamp.
%
After a deletion at level $\lambda$, every changed residue key is at least $\lambda$; therefore every resulting cascade at level $\lambda$ is processed before any live event of larger key can advance the level.
%
The queue execution consequently satisfies every premise of Theorem~\ref{thm:shell-order-replay}, which gives the from-scratch value of every residue pattern.
%
Theorem~\ref{thm:weighted-pattern-equivalence} transfers that value to all individual representatives.

Finally, each accepted nonstale event deletes one live pattern, and only finitely many stale events are generated per deletion because the pattern and leaf sets are finite.
%
Hence the queue empties and the algorithm terminates.
\end{proof}
